% problem_id: S001-lemma-cohomology-of-blowup
% source_id: S001 (Stacks Project)
% source_file: resolve.tex
% statement_locator: resolve.tex lines 568--582
% proof_locator: resolve.tex lines 584--698
% env: lemma
% pairing: tight (verified)
% status: complete (statement + author proof verbatim + reason + locators)
%
\begin{lemma}
\label{lemma-cohomology-of-blowup}
Let $(A, \mathfrak m, \kappa)$ be a regular local ring of dimension $2$.
Let $f : X \to S = \Spec(A)$ be the blowing up of $A$ in $\mathfrak m$.
Let $\mathcal{F}$ be a quasi-coherent $\mathcal{O}_X$-module.
\begin{enumerate}
\item $H^p(X, \mathcal{F}) = 0$ for $p \not \in \{0, 1\}$,
\item $H^1(X, \mathcal{O}_X(n)) = 0$ for $n \geq -1$,
\item $H^1(X, \mathcal{F}) = 0$ if $\mathcal{F}$ or $\mathcal{F}(1)$
is globally generated,
\item $H^0(X, \mathcal{O}_X(n)) = \mathfrak m^{\max(0, n)}$,
\item $\text{length}_A H^1(X, \mathcal{O}_X(n)) = -n(-n - 1)/2$
if $n < 0$.
\end{enumerate}
\end{lemma}

\begin{proof}
If $\mathfrak m = (x, y)$, then $X$ is covered by the spectra
of the affine blowup algebras $A[\frac{\mathfrak m}{x}]$ and
$A[\frac{\mathfrak m}{y}]$ because $x$ and $y$ placed in degree $1$
generate the Rees algebra $\bigoplus \mathfrak m^n$ over $A$.
See Divisors, Lemma \ref{divisors-lemma-blowing-up-affine} and
Constructions, Lemma \ref{constructions-lemma-proj-quasi-compact}.
Since $X$ is separated by
Constructions, Lemma \ref{constructions-lemma-proj-separated}
we see that cohomology of quasi-coherent sheaves vanishes in
degrees $\geq 2$ by Cohomology of Schemes, Lemma
\ref{coherent-lemma-vanishing-nr-affines}.

\medskip\noindent
Let $i : E \to X$ be the exceptional divisor, see
Divisors, Definition \ref{divisors-definition-blow-up}.
Recall that $\mathcal{O}_X(-E) = \mathcal{O}_X(1)$ is
$f$-relatively ample, see
Divisors, Lemma \ref{divisors-lemma-blowing-up-gives-effective-Cartier-divisor}.
Hence we know that $H^1(X, \mathcal{O}_X(-nE)) = 0$ for some $n > 0$,
see Cohomology of Schemes, Lemma \ref{coherent-lemma-kill-by-twisting}.
Consider the filtration
$$
\mathcal{O}_X(-nE) \subset \mathcal{O}_X(-(n - 1)E) \subset
\ldots \subset \mathcal{O}_X(-E) \subset \mathcal{O}_X \subset \mathcal{O}_X(E)
$$
The successive quotients are the sheaves
$$
\mathcal{O}_X(-t E)/\mathcal{O}_X(-(t + 1)E) =
\mathcal{O}_X(t)/\mathcal{I}(t) =
i_*\mathcal{O}_E(t)
$$
where $\mathcal{I} = \mathcal{O}_X(-E)$ is the ideal sheaf of $E$.
By Lemma \ref{lemma-blowup} we have $E = \mathbf{P}^1_\kappa$ and
$\mathcal{O}_E(1)$ indeed corresponds to the usual Serre twist of
the structure sheaf on $\mathbf{P}^1$. Hence the cohomology
of $\mathcal{O}_E(t)$ vanishes in degree $1$ for $t \geq -1$, see
Cohomology of Schemes, Lemma
\ref{coherent-lemma-cohomology-projective-space-over-ring}.
Since this is equal to $H^1(X, i_*\mathcal{O}_E(t))$ (by
Cohomology of Schemes, Lemma \ref{coherent-lemma-relative-affine-cohomology})
we find that $H^1(X, \mathcal{O}_X(-(t + 1)E)) \to H^1(X, \mathcal{O}_X(-tE))$
is surjective for $t \geq -1$. Hence
$$
0 = H^1(X, \mathcal{O}_X(-nE))
\longrightarrow
H^1(X, \mathcal{O}_X(-tE)) = H^1(X, \mathcal{O}_X(t))
$$
is surjective for $t \geq -1$ which proves (2).

\medskip\noindent
Let $\mathcal{F}$ be globally generated. This means there exists
a short exact sequence
$$
0 \to \mathcal{G} \to \bigoplus\nolimits_{i \in I} \mathcal{O}_X
\to \mathcal{F} \to 0
$$
Note that $H^1(X, \bigoplus_{i \in I} \mathcal{O}_X) =
\bigoplus_{i \in I} H^1(X, \mathcal{O}_X)$ by
Cohomology, Lemma \ref{cohomology-lemma-quasi-separated-cohomology-colimit}.
By part (2) we have $H^1(X, \mathcal{O}_X) = 0$.
If $\mathcal{F}(1)$ is globally generated, then we can find a
surjection $\bigoplus_{i \in I} \mathcal{O}_X(-1) \to \mathcal{F}$
and argue in a similar fashion.
In other words, part (3) follows from part (2).

\medskip\noindent
For part (4) we note that for all $n$ large enough we have
$\Gamma(X, \mathcal{O}_X(n)) = \mathfrak m^n$, see
Cohomology of Schemes, Lemma \ref{coherent-lemma-recover-tail-graded-module}.
If $n \geq 0$, then we can use the short exact sequence
$$
0 \to \mathcal{O}_X(n) \to \mathcal{O}_X(n - 1) \to
i_*\mathcal{O}_E(n - 1) \to 0
$$
and the vanishing of $H^1$ for the sheaf on the left to get a commutative
diagram
$$
\xymatrix{
0 \ar[r] &
\mathfrak m^{\max(0, n)} \ar[r] \ar[d] &
\mathfrak m^{\max(0, n - 1)} \ar[r] \ar[d] &
\mathfrak m^{\max(0, n)}/\mathfrak m^{\max(0, n - 1)} \ar[r] \ar[d] & 0\\
0 \ar[r] &
\Gamma(X, \mathcal{O}_X(n)) \ar[r] &
\Gamma(X, \mathcal{O}_X(n - 1)) \ar[r] &
\Gamma(E, \mathcal{O}_E(n - 1)) \ar[r] & 0
}
$$
with exact rows. In fact, the rows are exact also for $n < 0$
because in this case the groups on the right are zero.
In the proof of Lemma \ref{lemma-blowup}
we have seen that the right vertical arrow is an isomorphism
(details omitted). Hence if the left vertical arrow is an isomorphism, so
is the middle one. In this way we see that (4) holds by
descending induction on $n$.

\medskip\noindent
Finally, we prove (5) by descending induction on $n$ and the sequences
$$
0 \to \mathcal{O}_X(n) \to \mathcal{O}_X(n - 1) \to
i_*\mathcal{O}_E(n - 1) \to 0
$$
Namely, for $n \geq -1$ we already know $H^1(X, \mathcal{O}_X(n)) = 0$.
Since
$$
H^1(X, i_*\mathcal{O}_E(-2)) =
H^1(E, \mathcal{O}_E(-2)) =
H^1(\mathbf{P}^1_\kappa, \mathcal{O}(-2)) \cong \kappa
$$
by Cohomology of Schemes, Lemma
\ref{coherent-lemma-cohomology-projective-space-over-ring}
which has length $1$ as an $A$-module, we conclude from the long exact
cohomology sequence that (5) holds for $n = -2$. And so on and so forth.
\end{proof}
